\chapter{Thermodynamics of the Mixing Model}\label{app:thermo}

This appendix presents a self-contained theoretical framework for the hydrogen--silicate
mixing model used in \texttt{Atmodeller}. It unifies the free energy, the composition basis,
the chemical potentials and the stability conditions. The goal is to provide a physically
transparent explanation of why the binodal and the spinodal of the published model behave as
they do, and of which parts of that behaviour come from the model form rather than from the
data.

\section{Chemical Potentials on the Oxide Basis}

For a binary mixture of oxide units (component 1) and H$_2$ (component 2), with
$g_{\rm mix}(x)$ per mole of mixing units and $x$ the H$_2$ fraction, the chemical
potentials relative to the pure components are
\begin{equation}
\mu_1 = g_{\rm mix} - x\,g_{\rm mix}', \qquad \mu_2 = g_{\rm mix} + (1-x)\,g_{\rm mix}'.
\end{equation}
For Eq.~(\ref{eq:gs26}) with $f = 1-T/\tau+P/\pi$, the excess parts are
\begin{align}
\mu_1^{\rm ex} &= f\,x^2\left[B + 2(A-B)(1-x)\right],\\
\mu_2^{\rm ex} &= f\,(1-x)^2\left[A + 2(B-A)x\right],
\end{align}
which is equation~14 of \cite{gilmore2026} written for both components.

\section{Conversion to Native Formula Units}

\texttt{Atmodeller} counts $n_{\rm MgSiO_3}$ and $n_{\rm H_2}$. With
$G = (2n_{\rm MgSiO_3}+n_{\rm H_2})\,g_{\rm mix}(x)$ and
$x = n_{\rm H_2}/(2n_{\rm MgSiO_3}+n_{\rm H_2})$,
\begin{equation}
\frac{\partial G}{\partial n_{\rm MgSiO_3}} = 2\mu_1, \qquad
\frac{\partial G}{\partial n_{\rm H_2}} = \mu_2 .
\end{equation}
Per native mole, $h(y) = (2-y)\,g_{\rm mix}(x(y))$ with $x = y/(2-y)$ and
$dx/dy = 2/(2-y)^2$. Differentiating twice gives
\begin{equation}
h''(y) = \frac{4\,g_{\rm mix}''(x)}{(2-y)^3}.
\end{equation}
The curvature changes magnitude but not sign, so the spinodal is the same in both bases, as it
must be. The coexisting compositions are also the same once converted with
Eq.~(\ref{eq:basis}). Importing the published parameters without the factor $2-y$, as in our
incorrect-basis control, changes the free energy itself and therefore moves both curves.

\section{Spinodal Temperature and Its Threshold}

The excess curvature of Eq.~(\ref{eq:w}) is linear in composition,
\begin{equation}
W''(x) = 2B - 4A + 6(A-B)\,x .
\end{equation}
Setting $g_{\rm mix}'' = 0$ gives Eq.~(\ref{eq:tsp}). A positive spinodal temperature
requires $W''(x) < 0$, that is
\begin{equation}
x > x_{\rm min} = \frac{2B-4A}{6(B-A)} = 0.262
\end{equation}
for the published parameters. This corresponds to $w = \pubThreshold{}$\,wt\% local hydrogen.
Below this threshold a parcel can still separate in equilibrium, but its homogeneous state
never becomes unstable: separation there always requires nucleation.

\section{Why the Withheld Fraction Does Not Depend on Pressure}

Write $g_{\rm mix}/(RT) = s(x) + r\,W(x)$, with $s$ the ideal term and
\begin{equation}
r(T,P) = \frac{1 - T/\tau + P/\pi}{RT}.
\end{equation}
Phase equilibrium and stability depend on $T$ and $P$ only through $r$. At fixed composition
the spinodal fixes $r_{\rm sp} = -1/[x(1-x)W''(x)]$, and the binodal fixes another value
$r_b(x)$. Solving $r(T,P) = r_\ast$ gives
\begin{equation}
T_\ast(x,P) = \frac{1+P/\pi}{R\,r_\ast + 1/\tau} = T_\ast(x,0)\left(1+\frac{P}{\pi}\right).
\end{equation}
Both onset temperatures therefore scale with $1+P/\pi$, and the coexisting compositions at the
spinodal, hence $f_{\rm H}$, are the same at every pressure. The numerical invariance of
$f_{\rm H}$ (Section~\ref{sec:inventory}) is a consequence of this separable form, not
independent evidence. Classes M3 and M4 share it, because all their terms scale with the
same factor; class M2 does not.

\section{Critical Point}

At the critical point both the second and third derivatives of $g_{\rm mix}$ vanish. For
Eq.~(\ref{eq:gs26}) this gives equations~17--19 of \cite{gilmore2026}: the ratio $B/A$ fixes
the critical composition, here $x_C = \pubXc{}$ ($w = \pubWc{}$\,wt\%), and
\begin{equation}
T_C(P) = T_C^0\left(1+\frac{P}{\pi}\right), \qquad T_C^0 = \pubTcZero{}\,{\rm K}.
\end{equation}
The critical temperature is the same whichever component carries $A$; the critical
composition is not.

\section{Curvature Change of the Identifiability Perturbation}

For the perturbation of Eq.~(\ref{eq:bump}) with gap width $\Delta x = x_{\rm rich} -
x_{\rm poor}$,
\begin{equation}
\frac{\partial^2 \delta g}{\partial x^2}
= \lambda B_{\rm gap}\,\frac{384\,u(1-u)(1-5u+5u^2)}{\Delta x^2}.
\end{equation}
The change is negative in the middle of the gap ($0.28 < u < 0.72$) and positive near its
edges, so the perturbation can move the spinodal in either direction. This is why the 3\,wt\% spinodal
rises while the 5\,wt\% spinodal falls in Figure~\ref{fig:identifiability}. The value,
first and second derivatives of $\delta g$ vanish at $u=0$ and $u=1$, which keeps every
equilibrium quantity unchanged.
